\expchapter{实验二\quad 中国地图四染色问题}
\sect{教学目标}
\noindent\hspace*{2em}知识和能力目标：熟练掌握和应用栈结构解决实际问题，本例主要是从递归算法到非递归算法的问题。\par
\noindent\hspace*{2em}价值目标：\par
\noindent\hspace*{2em}科学精神：希望你了解四色定理的产生及研究历程，理解科学研究的严谨性、长期性及其对人类社会发展的深远影响。\par
\noindent\hspace*{2em}人文素质：你是否知道什么是飞地？\par
\sect{问题描述}
\noindent\hspace*{2em}四色问题又称四色猜想、四色定理，是世界近代三大数学难题之一。1852 年毕业于伦敦大学的佛朗西斯·格斯里（Francis Guthrie）在一家科研单位进行地图着色工作时，发现了一个有趣现象：每幅地图都可以只用四种颜色着色，就能使有共同边界的国家都着上不同颜色（不含飞地\footnote{飞地，是一种特殊的人文地理现象，指隶属于某一行政区管辖但不与本区毗连的土地。}）。那么，这是一个因地图不够复杂的巧合？还是一个确认的定理呢？\par
\noindent\hspace*{2em}多年以来，数学家们为证明四色定理绞尽脑汁，所引进的概念与方法刺激了拓扑学与图论的发展。在四色问题研究中，不少新的数学理论和数学计算技巧随之产生，例如将地图的着色问题转化为图论问题，丰富了图论的内容。不仅如此，"四色问题"在有效设计航空班机日程表、设计计算机的编码程序上都起到了推动作用。\par
\noindent\hspace*{2em}中华人民共和国有 34 个省级行政区，含 23 个省、5 个自治区、4 个直辖市、2 个特别行政区\footnote{中国行政区划（2020 年），行政区划网。}。\par
\begin{center}
\includegraphics[width=0.65\textwidth]{images/exp2-1.png}
\end{center}
\noindent\hspace*{2em}23 个省分别为：河北省、山西省、辽宁省、吉林省、黑龙江省、江苏省、浙江省、安徽省、福建省、江西省、山东省、河南省、湖北省、湖南省、广东省、海南省、四川省、贵州省、云南省、陕西省、甘肃省、青海省、台湾省。\par
\noindent\hspace*{2em}5 个自治区分别为：内蒙古自治区、广西壮族自治区、西藏自治区、宁夏回族自治区、新疆维吾尔自治区。\par
\noindent\hspace*{2em}4 个直辖市分别为：北京市、天津市、上海市、重庆市。\par
\noindent\hspace*{2em}2 个特别行政区分别为：香港特别行政区、澳门特别行政区。本问题的任务是：编写程序，用四种颜色对中国地图染色，输出任意一种染色方案。\par
\noindent\hspace*{2em}是否能用 3 种颜色对中国地图染色？请编程验证。假设如下：\par
\noindent\hspace*{2em}\textbf{四种颜色不限，输出形式不限}\par
\noindent\hspace*{2em}\textbf{用递归和非递归两种算法来实现}\par
\sect{基本思路}
\noindent\hspace*{2em}地图可用图的邻接矩阵方法来存储。染色过程可借助一个栈结构，由于省级行政区数已知为 34，因此可使用顺序栈。\par
\noindent\hspace*{2em}算法步骤（以递归版本为例）：\par
\noindent\hspace*{2em}Step1：对区域下标 $i$ 从 0 开始依次尝试。\par
\noindent\hspace*{2em}Step2：依次试探四种颜色 1、2、3、4，只要与所有已着色邻居不冲突即可着上。\par
\noindent\hspace*{2em}Step3：若当前区域 1\~{}4 号色均冲突，则回溯到上一区域重新选择颜色。\par
\noindent\hspace*{2em}Step4：若所有区域都被成功着色，则得到一个合法方案；否则无解。\par
\noindent\hspace*{2em}非递归版本使用顺序栈保存"当前正在尝试着色的区域下标"以及"该区域下一个待试探的颜色"，遇冲突时通过出栈实现回溯。\par
\noindent\hspace*{2em}染色时不妨总是从 1 号色开始试探。\par
\sect{实验结果及分析}
\noindent\hspace*{2em}1. 实验结果\par
\noindent\hspace*{2em}（1）实验过程描述\par
\noindent\hspace*{2em}1. 本实验用 Rust 语言开发，使用 Cargo 构建系统和 rustc 编译器，运行机器为本人笔记本电脑（macOS，Apple Silicon）。本人独立完成编码和测试，参考相关算法资料关于递归回溯的描述确定了方案。\par
\noindent\hspace*{2em}2. 测试数据选取思路为：典型、覆盖边界情况、覆盖算法变体。共选取 3 组数据进行测试：(a) 34 省级行政区完整邻接矩阵的中国地图，验证 4 着色可行性；(b) 构造的 K4 完全图（4 个节点两两相邻），验证 4 色必要性；(c) 3 着色尝试，验证中国地图不能用 3 种颜色着色。\par
\noindent\hspace*{2em}3. 4 染色回溯和 3 染色可行性验证的逻辑是测试重点。因为 4 染色和 3 染色共用同一回溯框架，区别仅在试探的颜色数。\par
\noindent\hspace*{2em}（2）测试结果\par
\noindent\hspace*{2em}运行实验程序对中国地图进行 4 染色和 3 染色尝试，结果如下表：\par
\begin{center}
\fontsize{9pt}{12pt}\selectfont
\renewcommand{\arraystretch}{1.25}
\begin{tabular}{|c|c|c|c|}
\hline
测试项 & 输入 & 输出 & 结论 \\ \hline
递归 4 着色 & 中国 34 行政区 & 合法方案（34 项） & 成功 \\ \hline
非递归 4 着色 & 中国 34 行政区 & 合法方案（34 项） & 成功 \\ \hline
3 着色 & 中国 34 行政区 & 无解 & 不可行 \\ \hline
\end{tabular}
\end{center}
\noindent\hspace*{2em}两种算法均得到一个合法方案（颜色取值仅为 1\~{}4），3 着色搜索失败——证明中国地图不能用 3 种颜色染色。递归与非递归 4 着色方案如下：\par
\begin{center}
\fontsize{8.5pt}{11pt}\selectfont
\renewcommand{\arraystretch}{1.15}
\begin{tabular}{|c|c|c||c|c|c|}
\hline
\textbf{行政区} & \textbf{递归色号} & \textbf{非递归色号} & \textbf{行政区} & \textbf{递归色号} & \textbf{非递归色号} \\ \hline
北京   & 1 & 1 & 湖北   & 3 & 3 \\ \hline
天津   & 2 & 2 & 湖南   & 1 & 1 \\ \hline
河北   & 3 & 3 & 广东   & 3 & 3 \\ \hline
山西   & 1 & 1 & 广西   & 2 & 2 \\ \hline
内蒙古 & 2 & 2 & 海南   & 1 & 1 \\ \hline
辽宁   & 1 & 1 & 重庆   & 2 & 2 \\ \hline
吉林   & 3 & 3 & 四川   & 1 & 1 \\ \hline
黑龙江 & 1 & 1 & 贵州   & 3 & 3 \\ \hline
上海   & 1 & 1 & 云南   & 4 & 4 \\ \hline
江苏   & 2 & 2 & 西藏   & 2 & 2 \\ \hline
浙江   & 3 & 3 & 陕西   & 4 & 4 \\ \hline
安徽   & 1 & 1 & 甘肃   & 3 & 3 \\ \hline
福建   & 1 & 1 & 青海   & 4 & 4 \\ \hline
江西   & 2 & 2 & 宁夏   & 1 & 1 \\ \hline
山东   & 4 & 4 & 新疆   & 1 & 1 \\ \hline
河南   & 2 & 2 & 台湾   & 1 & 1 \\ \hline
 &   &   & 香港   & 1 & 1 \\ \hline
 &   &   & 澳门   & 1 & 1 \\ \hline
\end{tabular}
\end{center}
\noindent\hspace*{2em}对得到的 4 着色方案进行合法性校验（检查任意一对有共同边界的省份是否异色）以及颜色取值范围校验（仅取 1\~{}4），均通过。3 着色在递归版本上搜索完所有可能后返回无解，非递归版本同样返回 None。\par
\noindent\hspace*{2em}将 4 染色方案以区域编号（0\~{}33）为横轴、颜色号为纵轴绘制分组柱状图：\par
\begin{center}
\begin{tikzpicture}
\begin{axis}[
    ybar,
    width=0.98\textwidth,
    height=5.8cm,
    ymin=0, ymax=4.5,
    ytick={1,2,3,4},
    ylabel={颜色号},
    xlabel={区域编号},
    symbolic x coords={0,1,2,3,4,5,6,7,8,9,10,11,12,13,14,15,16,17,18,19,20,21,22,23,24,25,26,27,28,29,30,31,32,33},
    xtick={0,5,10,15,20,25,30,33},
    bar width=2.2pt,
    enlarge x limits=0.04,
    grid=major,
    legend pos=outer north east,
    legend style={font=\fontsize{8pt}{10pt}\selectfont},
]
\addplot[fill=blue!60] coordinates {
    (0,1)(1,2)(2,3)(3,1)(4,2)(5,1)(6,3)(7,1)(8,1)(9,2)(10,3)(11,1)(12,1)(13,2)(14,4)(15,2)(16,3)(17,1)(18,3)(19,2)(20,1)(21,2)(22,1)(23,3)(24,4)(25,2)(26,4)(27,3)(28,4)(29,1)(30,1)(31,1)(32,1)(33,1)
};
\addlegendentry{递归 4 着色}
\addplot[fill=red!60] coordinates {
    (0,1)(1,2)(2,3)(3,1)(4,2)(5,1)(6,3)(7,1)(8,1)(9,2)(10,3)(11,1)(12,1)(13,2)(14,4)(15,2)(16,3)(17,1)(18,3)(19,2)(20,1)(21,2)(22,1)(23,3)(24,4)(25,2)(26,4)(27,3)(28,4)(29,1)(30,1)(31,1)(32,1)(33,1)
};
\addlegendentry{非递归 4 着色}
\end{axis}
\end{tikzpicture}
\end{center}
\noindent\hspace*{2em}从分组柱状图可见，两种算法在 34 个区域上得到完全相同的着色方案（蓝红两柱等高并列），所有区域颜色号均落在 1\~{}4 之间。下表统计了各颜色被使用的次数：\par
\begin{center}
\fontsize{9pt}{12pt}\selectfont
\renewcommand{\arraystretch}{1.25}
\begin{tabular}{|c|c|c|c|c|}
\hline
\textbf{颜色号} & \textbf{1} & \textbf{2} & \textbf{3} & \textbf{4} \\ \hline
使用区域数 & 15 & 8 & 7 & 4 \\ \hline
\end{tabular}
\end{center}
\noindent\hspace*{2em}由统计可见，1 号色被使用 15 次最多，2 号色 8 次次之，4 号色仅 4 次最少。这是因为递归/非递归算法均采用固定下标顺序（0\~{}33）依次试探，从 1 号色起着色，故位于"着色序列前段"的区域有更多机会取到 1 号色。\par
\noindent\hspace*{2em}2. 结果及方案分析\par
\noindent\hspace*{2em}（1）本实验的重点及难点\par
\noindent\hspace*{2em}核心工作是建立 34$\times$34 邻接矩阵描述 34 个省级行政区的相邻关系，按编号顺序依次试探颜色，冲突时回溯到上一区域。\par
\noindent\hspace*{2em}非递归改写中需要用顺序栈显式保存"待回溯的中间状态"——栈中记录"当前区域下标"和"下一个待试探的颜色号"，冲突时出栈实现状态回退。\par
\noindent\hspace*{2em}（2）本实验设计的优点及可以改进的方向\par
\noindent\hspace*{2em}程序支持 4 染色和 3 染色验证。(1) 邻接矩阵与区域名称集中存储在常量数组中；(2) 递归与非递归共用 is\_safe 检测函数；(3) 单元测试覆盖 K4 完全图与零边图；(4) 非递归版本显式使用顺序栈。\par
\noindent\hspace*{2em}设计的可优化处包括：(1) 当前按固定下标顺序着色，回溯次数可能较多，可改用 Welsh-Powell 等启发式策略按度数从大到小选择下一个区域着色，平均搜索量更小；(2) 邻接矩阵存储为 34$\times$34=1156 字节，对中国地图规模可以接受，但若推广到更细粒度的县级行政区，可改用邻接表节省空间；(3) 非递归版本代码较为冗长，栈操作和回溯控制逻辑可以进一步抽象。\par
\noindent\hspace*{2em}3. 从四色定理看科学研究方法。\par
\noindent\hspace*{2em}四色猜想从 1852 年提出到 1976 年 Appel 和 Haken 用计算机验证约 1936 个构型完成证明，历时 124 年。计算机辅助证明已成为数学研究的新范式。四色理论已用于 PCB 布线、航班排程和无线频率分配。\par
\noindent\hspace*{2em}作为计算机专业的学生，应重视算法理论基础，也关注工程落地——如图着色之于排课系统。面对复杂问题时，可借助工具处理繁琐环节，将精力集中在创造性思路上。\par
\noindent\hspace*{2em}4. 附件：源程序代码\par
\noindent\hspace*{2em}4. 附件：源程序代码\par
\noindent\hspace*{2em}实验代码使用 Rust 语言编写，Cargo 构建。整个工程分为三个源文件：lib.rs 负责核心算法，stack.rs 负责顺序栈抽象，main.rs 负责命令行入口。\par
{
\fontsize{7.5pt}{10pt}\selectfont
\setlength{\parindent}{0pt}
\setlength{\columnsep}{0.6cm}
\begin{multicols}{2}
\noindent\textbf{src/lib.rs}\par

\begin{verbatim}
pub mod region;
pub mod stack;

use crate::region::{adjacency, N};

/// 检查颜色 color[idx] = c 是否与所有
/// 已着色的邻居冲突
fn is_safe(map: &[[u8; N]; N],
           color: &[u8; N],
           idx: usize, c: u8) -> bool {
    for j in 0..N {
        if map[idx][j] == 1
           && color[j] != 0
           && color[j] == c {
            return false;
        }
    }
    true
}

/// 递归回溯求解 n-着色问题
pub fn color_recursive(
        map: &[[u8; N]; N],
        n_colors: u8
    ) -> Option<[u8; N]> {
    let mut color = [0u8; N];

    if try_color(map, &mut color, 0,
                 n_colors) {
        Some(color)
    } else {
        None
    }
}

fn try_color(map: &[[u8; N]; N],
             color: &mut [u8; N],
             idx: usize,
             n_colors: u8) -> bool {
    if idx == N {
        return true;
    }

    for c in 1..=n_colors {
        if is_safe(map, color, idx, c) {
            color[idx] = c;
            if try_color(map, color,
                         idx + 1,
                         n_colors) {
                return true;
            }
            color[idx] = 0;
        }
    }

    false
}

/// 非递归回溯：借助顺序栈保存
/// 当前正在处理的区域下标
pub fn color_non_recursive(
        map: &[[u8; N]; N],
        n_colors: u8
    ) -> Option<[u8; N]> {
    let mut color = [0u8; N];
    let mut st =
        stack::Stack::new(N + 1);

    let mut idx: usize = 0;
    loop {
        // 找一个尚未着色的区域
        while idx < N
              && color[idx] != 0 {
            idx += 1;
        }

        if idx == N {
            return Some(color);
        }

        // 从 color[idx]+1 开始尝试
        let mut c = color[idx] + 1;
        let mut found = false;

        while c <= n_colors {
            if is_safe(map, &color,
                       idx, c) {
                found = true;
                break;
            }
            c += 1;
        }

        if found {
            color[idx] = c;
            st.push(idx);
            idx += 1;
        } else {
            // 回溯
            color[idx] = 0;

            let mut bt_ok = false;
            while let Some(prev) = st.pop() {
                let prev_color = color[prev];
                color[prev] = 0;
                let mut nc = prev_color + 1;
                let mut ok = false;
                while nc <= n_colors {
                    if is_safe(map, &color,
                               prev, nc) {
                        ok = true;
                        break;
                    }
                    nc += 1;
                }

                if ok {
                    color[prev] = nc;
                    st.push(prev);
                    idx = prev + 1;
                    bt_ok = true;
                    break;
                }
            }

            if !bt_ok {
                return None;
            }
        }
    }
}

/// 验证一个 n-着色方案是否合法
pub fn is_valid_coloring(
        map: &[[u8; N]; N],
        color: &[u8; N]) -> bool {
    for i in 0..N {
        for j in (i + 1)..N {
            if map[i][j] == 1
               && color[i] == color[j] {
                return false;
            }
        }
    }
    true
}

pub fn default_map() -> [[u8; N]; N] {
    adjacency()
}

pub use region::region_name;
\end{verbatim}

\columnbreak

\noindent\textbf{src/stack.rs}\par

\begin{verbatim}
/// 顺序栈，存储当前正在尝试
/// 着色的区域下标
pub struct Stack {
    data: Vec<usize>,
    top: usize,
}

impl Stack {
    pub fn new(capacity: usize)
        -> Self {
        Stack {
            data: vec![0; capacity],
            top: 0,
        }
    }

    pub fn push(&mut self,
                val: usize) -> bool {
        if self.top >= self.data.len() {
            return false;
        }
        self.data[self.top] = val;
        self.top += 1;
        true
    }

    pub fn pop(&mut self)
        -> Option<usize> {
        if self.top == 0 {
            return None;
        }
        self.top -= 1;
        Some(self.data[self.top])
    }

    pub fn peek(&self)
        -> Option<usize> {
        if self.top == 0 {
            None
        } else {
            Some(self.data[
                self.top - 1])
        }
    }

    pub fn is_empty(&self) -> bool {
        self.top == 0
    }

    pub fn len(&self) -> usize {
        self.top
    }
}
\end{verbatim}

\noindent\textbf{src/region.rs}\par

\begin{verbatim}
pub const REGIONS: [&str; 34] = [
    "北京", "天津", "河北", "山西",
    "内蒙古",
    "辽宁", "吉林", "黑龙江",
    "上海", "江苏", "浙江", "安徽",
    "福建", "江西", "山东",
    "河南", "湖北", "湖南", "广东",
    "广西", "海南",
    "重庆", "四川", "贵州", "云南",
    "西藏",
    "陕西", "甘肃", "青海", "宁夏",
    "新疆",
    "台湾", "香港", "澳门",
];

pub const N: usize = 34;

pub fn adjacency() -> [[u8; N]; N] {
    let mut m = [[0u8; N]; N];
    // 共 70 条无向边，省略具体列表
    // ...
    m
}

pub fn region_name(idx: usize)
    -> &'static str {
    REGIONS[idx]
}
\end{verbatim}

\noindent\textbf{src/main.rs}\par

\begin{verbatim}
use exp_2::{
    color_non_recursive,
    color_recursive,
    default_map,
    is_valid_coloring,
    region_name
};

fn print_coloring(
        label: &str,
        color: &[u8; 34]) {
    println!("\n=== {} ===", label);
    for i in 0..34 {
        println!(
            "  {:>4} -> 色 {}",
            region_name(i),
            color[i]
        );
    }
}

fn main() {
    let map = default_map();

    // 递归 4 着色
    if let Some(c) =
        color_recursive(&map, 4) {
        assert!(
            is_valid_coloring(&map, &c));
        print_coloring(
            "递归 4 着色方案", &c);
    }

    // 非递归 4 着色
    if let Some(c) =
        color_non_recursive(&map, 4) {
        assert!(
            is_valid_coloring(&map, &c));
        print_coloring(
            "非递归 4 着色方案", &c);
    }

    // 3 着色可行性验证
    let three =
        color_recursive(&map, 3);
    match three {
        Some(_) => println!(
            "3 着色：存在合法方案"),
        None => println!(
            "3 着色：无合法方案"),
    }
}
\end{verbatim}

\end{multicols}
}
\vspace{0.5cm}
{
\fontsize{7.5pt}{10pt}\selectfont
\setlength{\parindent}{0pt}
\setlength{\columnsep}{0.6cm}
\begin{multicols}{2}
\noindent\textbf{tests/test\_fn.rs}\par

\begin{verbatim}
use exp_2::{
    color_non_recursive,
    color_recursive,
    default_map,
    is_valid_coloring,
    region_name
};

fn k4_map() -> [[u8; 34]; 34] {
    let mut m = [[0u8; 34]; 34];
    let nodes = [0, 1, 2, 3];
    for &i in &nodes {
        for &j in &nodes {
            if i != j {
                m[i][j] = 1;
            }
        }
    }
    m
}

fn empty_map() -> [[u8; 34]; 34] {
    [[0u8; 34]; 34]
}

#[test]
fn recursive_4color_finds() {
    let map = default_map();
    let c = color_recursive(
        &map, 4
    ).expect("应可解");
    assert!(
        is_valid_coloring(&map, &c));
    for i in 0..34 {
        assert!((1..=4)
            .contains(&c[i]));
    }
}

#[test]
fn non_recursive_matches() {
    let map = default_map();
    let c1 = color_recursive(
        &map, 4).unwrap();
    let c2 = color_non_recursive(
        &map, 4).unwrap();
    assert!(
        is_valid_coloring(&map, &c1));
    assert!(
        is_valid_coloring(&map, &c2));
}

#[test]
fn three_color_should_fail() {
    let map = default_map();
    assert!(color_recursive(
        &map, 3).is_none());
    assert!(color_non_recursive(
        &map, 3).is_none());
}

#[test]
fn k4_needs_four() {
    let map = k4_map();
    assert!(color_recursive(
        &map, 3).is_none());
    let c = color_recursive(
        &map, 4).unwrap();
    assert!(
        is_valid_coloring(&map, &c));
}

#[test]
fn empty_map_one_color() {
    let map = empty_map();
    let c = color_recursive(
        &map, 1).unwrap();
    for v in c.iter() {
        assert_eq!(*v, 1);
    }
    assert!(
        is_valid_coloring(&map, &c));
}

#[test]
fn region_names_count() {
    let map = default_map();
    let c = color_recursive(
        &map, 4).unwrap();
    for i in 0..34 {
        assert!(
            !region_name(i).is_empty());
        assert!(c[i] >= 1);
    }
}

#[test]
fn non_border_pairs() {
    let map = default_map();
    let pairs = [
        (1,14), (3,14), (4,30),
        (12,31), (18,20), (19,20),
    ];
    for &(i, j) in pairs.iter() {
        assert_eq!(map[i][j], 0);
        assert_eq!(map[j][i], 0);
    }
}
\end{verbatim}

\end{multicols}
}
{
\fontsize{7.5pt}{10pt}\selectfont
\setlength{\parindent}{0pt}
\noindent\textbf{测试用例}\par
\noindent (a) 中国地图 34 行政区完整邻接矩阵：递归/非递归 4 着色均成功，3 着色无解。\par
\noindent (b) K4 完全图（4 节点两两相邻）：4 着色成功，3 着色无解。\par
\noindent (c) 零边图：1 着色即可（全部着 1 号色），无冲突。\par
\noindent (d) 非共同边界检查：天津-山东、山西-山东、内蒙古-新疆、福建-台湾、广东-海南、广西-海南均不记为相邻。\par
}

